% ============================ 6. Discussion =============================
\section{Discussion}\label{sec:discussion}

\subsection{What scalar temporal attention does compute}
\label{sec:disc-what}

Taken together, the frozen results characterize TAN as
\begin{equation*}
  \boxed{\text{prediction-error gating} \;+\;
    \text{history dependence} \;+\;
    \text{scalar temporal saliency weighting}.}
\end{equation*}
The state-collision experiments establish history dependence and show it
to be conditional on surprise: in gate-closed episodes the dynamics
collapse to exact LIF contraction, and in gate-open episodes the
windowed context is re-coupled into the membrane through an
attention-weighted current.  The geometry experiments show that this
re-coupling has an observable, event-locked signature in the effective
rank and covariance scale of the response representation, and that
adaptive attention adds structure beyond uniform temporal integration in
the frame that contains the attention coordinates.  The retrieval
experiment provides constrained evidence that such structure can support
target-relevant readout in content-rich windows, relative to the
no-attention control, in a task whose fixed-position layout
structurally favours a passive delay line (B2).

\subsection{What scalar temporal attention does not compute}
\label{sec:disc-whatnot}

The same body of evidence shows what the minimal kernel does not
provide.  In the shared observable coordinates, TAN does not possess a
higher effective rank of event responses than LIF; effective-rank
expansion is generic to history-carrying units.  On a fixed-position
retrieval task, a passive delay line outperforms attention variants.
And in the identifiability audit, the scalar kernel failed every
routing gate: attention focus did not follow the query, the
target/distractor contrast was inverted, and holding history fixed while
switching the query never changed which history tap was attended.  The
architecture therefore does not implement content-addressable
$Q$--$K$ routing, and semantic binding is not established by any result
in this paper.

\subsection{Why attention weighting is not routing}\label{sec:disc-why}

The mechanistic reason is the factorization of Section~\ref{sec:model-scalar}:
with a scalar surprise $S_t$, attention logits are a common scalar times
the window values, so their \emph{relative ordering} is a property of the
input amplitudes alone.  A query that enters only through $S_t$ can
modulate how sharply the distribution concentrates -- a gain-like,
saliency-like operation -- but cannot determine \emph{which} historical
item is selected.  Weighting strength and content selection are
therefore different operations, consistent with the interpretive
literature on attention weights in general \citep{jainwallace2019,
serranosmith2019,wiegreffe2019}: observed weights can be read as
explanations or routing decisions only if the mechanism is shown to
change selected content under a control signal.

\subsection{Representation complexity is not computational selectivity}
\label{sec:disc-complexity}

A recurrent theme across phases is that richer state structure does not
by itself confer selective computation.  Phase~2 showed representational
reorganization in attention-bearing coordinates; Probe~1 showed that a
fixed memory baseline can outperform attention on a storage-aligned
task; Probe~2 showed that a scalar weighting mechanism with nontrivial
history dependence cannot select content by query.  These results
caution against inferring computational capability from state or
representation complexity alone \citep{chungabbott2021,gaoganguli2015},
and against treating dimensionality-like statistics as intrinsic
capacity measures \citep{facco2017}.  We make no complexity theorem here;
the empirical point is that the layers of the hierarchy in
Section~\ref{sec:intro} must each be verified directly.

\subsection{Relation to dynamical systems}\label{sec:disc-dyn}

From a dynamical-systems perspective, TAN is a non-autonomous,
input-driven map whose contraction is conditional on a prediction-error
gate: leaky decay guarantees boundedness \citep{lohmiller1998,
strogatz2014}, while windowed updates make the effective state
higher-dimensional and the one-step divergence of near-colliding states
reflect the local expansion of the driven map.  Reset re-bifurcation
shows that the spiking readout does not erase the windowed memory,
connecting the single-neuron setting to delay- and
attractor-based memory accounts \citep{compte2000,durstewitz2000,
elman1990}.  Standard nonlinear time-series tools (embedding, effective
dimension) \citep{takens1981,grassberger1983,kennel1992,facco2017,
bandt2002} motivated our observable choices, but we report only
estimator-level statistics (effective rank, log-trace, spectra) with the
caveats stated in Section~\ref{sec:limitations}.

\subsection{Implications for multi-neuron architectures}\label{sec:disc-org}

The mechanistic boundary of the scalar kernel identifies what an
organizing principle would need to add: a mechanism in which a query can
change the \emph{relative} compatibility of candidate memories.
Candidate organizational directions -- distinct from the present results
-- include multiple dendritic compartments \citep{london2005,larkum2013,
spruston2008}, opponent/excitatory--inhibitory circuits, lateral
inhibition and competitive microcircuits, and phase- or synchrony-based
binding \citep{von_der_malsburg1995,singer1995,fries2005}.  All such
directions are future hypotheses; none is tested or claimed here.
